\documentclass[a4paper, 12pt]{article}

\usepackage{amsmath, amsthm, amssymb}
\usepackage{unicode-math}
\usepackage{xltxtra}
\usepackage{xgreek}
\setmainfont{Calibri}

\usepackage{fancyhdr}
\renewcommand{\footrulewidth}{0.4pt}
\usepackage{lastpage}

\usepackage[a4paper, top=2.5cm, bottom=2.5cm, left=2.5cm, right=2.5cm, headheight=1cm, footskip=1.5cm]{geometry}

\usepackage{tabularx}
\usepackage{enumitem}

% Ρυθμίσεις Επικεφαλίδων
\pagestyle{fancy}
\fancyhf{}
\rhead{Άλγεβρα Α' Λυκείου}
\lhead{Γραφική Παράσταση Συνάρτησης}
\cfoot{Σελίδα \thepage \ από \pageref{LastPage}}

\begin{document}

% ----------------------------------------------------------------------
% 1. ΣΧΕΔΙΟ ΜΑΘΗΜΑΤΟΣ (Lesson Plan)
% ----------------------------------------------------------------------
\begin{center}
  \LARGE \textbf{Σχέδιο Μαθήματος} \\
  \large \textbf{Ενότητα: Η Γραφική Παράσταση μιας Συνάρτησης ($C_f$)} \\
  \vspace{0.2cm}
  \small \textit{Διάρκεια: 45 λεπτά}
\end{center}

\section*{Στόχοι Μαθήματος}
\begin{enumerate}
  \item \textbf{Εννοιολογική Σύνδεση}: Κατανόηση ότι η $C_f$ αποτελείται από όλα τα σημεία $(x, y)$ που ικανοποιούν τη σχέση $y=f(x)$.
  \item \textbf{Οπτική Αντίληψη}: Αναγνώριση του Πεδίου Ορισμού και του Συνόλου Τιμών πάνω στους άξονες.
  \item \textbf{Γραφική Επίλυση}: Ικανότητα προσδιορισμού ριζών και προσήμου της συνάρτησης απευθείας από το γράφημα.
  \item \textbf{Συγκριτική Ανάγνωση}: Εντοπισμός των σημείων τομής δύο γραφικών παραστάσεων και ερμηνεία τους.
\end{enumerate}

\section*{Ροή Μαθήματος (Χρονοδιάγραμμα)}
\begin{itemize}[label=--]
  \item \textbf{00'--07' (Διερεύνηση):} Προβολή ενός κενού συστήματος αξόνων. Ερώτημα: «Αν είχαμε τη δυνατότητα να σχεδιάσουμε χιλιάδες σημεία $M(x, f(x))$ σε ένα δευτερόλεπτο, τι θα βλέπαμε;». Σύντομη επίδειξη του Excel μόνο για την ταχεία παραγωγή σημείων που σχηματίζουν μια καμπύλη.
  \item \textbf{07'--15' (Θεωρητική Πλαισίωση):} Ορισμός $C_f$. Συζήτηση για τη σημασία της τεταγμένης: «Το ύψος κάθε σημείου είναι η τιμή της συνάρτησης».
  \item \textbf{15'--35' (Εργασία σε Ομάδες):} Επεξεργασία του φύλλου εργασίας. Εστίαση στην ανάγνωση των έτοιμων σχημάτων από τις ασκήσεις της παρουσίασης (Ασκήσεις 1, 3, 5 και 6).
  \item \textbf{35'--40' (Σύνοψη):} Συζήτηση για το πώς το γράφημα μας επιτρέπει να «λύνουμε» εξισώσεις χωρίς να κάνουμε πράξεις.
  \item \textbf{40'--45' (Ατομική Αξιολόγηση):} Συμπλήρωση του Exit Ticket για τον έλεγχο κατανόησης των εννοιών.
\end{itemize}

\newpage

% ----------------------------------------------------------------------
% 2. ΦΥΛΛΟ ΕΡΓΑΣΙΑΣ (Worksheet)
% ----------------------------------------------------------------------
\begin{center}
  \Large \textbf{Φύλλο Εργασίας: Ανακαλύπτοντας τη $C_f$} \\
  \small Τμήμα: \dots\dots \quad Ομάδα: \dots\dots\dots\dots\dots\dots
\end{center}

\vspace{0.5cm}

\textbf{Δραστηριότητα 1: Η γεωμετρική ερμηνεία του τύπου} \\
Δίνεται μια συνάρτηση $f$. Χρησιμοποιώντας τη λογική «το σημείο $M(x,y)$ ανήκει στη $C_f$ αν και μόνο αν $y=f(x)$»:
\begin{enumerate}
  \item Αν γνωρίζετε ότι $f(2)=5$, ποιο σημείο μπορείτε να σχεδιάσετε σίγουρα στο επίπεδο; \dots\dots\dots
  \item Αν δείτε στο γράφημα ότι η καμπύλη περνάει από το σημείο $(-1, 4)$, τι συμπέρασμα βγάζετε για την τιμή $f(-1)$; \dots\dots\dots\dots\dots\dots
\end{enumerate}

\vspace{0.8cm}

\textbf{Δραστηριότητα 2: Ανάγνωση Πληροφοριών (Ασκήσεις 1 \& 3)} \\
Παρατηρήστε τις γραφικές παραστάσεις στις αντίστοιχες ασκήσεις της παρουσίασης:
\begin{enumerate}
  \item Ποια «σκιά - προβολή» αφήνει η $C_f$ πάνω στον άξονα $x'x$; Τι αντιπροσωπεύει αυτή η «σκιά» για τη συνάρτηση; \\
        \textit{Απάντηση:} \dots\dots\dots\dots\dots\dots\dots\dots
  \item Εντοπίστε τα σημεία τομής με τον $x'x$. Τι τιμή έχει η συνάρτηση σε αυτά τα σημεία; \\
        \textit{Απάντηση:} \dots\dots\dots\dots\dots\dots\dots\dots
\end{enumerate}

\vspace{0.8cm}

\textbf{Δραστηριότητα 3: Σύγκριση και Πρόσημο (Ασκήσεις 5 \& 6)} \\
Κοιτάζοντας τα σχήματα όπου εμφανίζονται δύο συναρτήσεις $f$ και $g$:
\begin{enumerate}
  \item Πώς ερμηνεύετε γραφικά την πρόταση «το $x_0$ είναι λύση της εξίσωσης $f(x)=g(x)$»; \\
        \textit{Απάντηση:} \dots\dots\dots\dots\dots\dots\dots\dots
  \item Πώς καταλαβαίνουμε από το σχήμα για ποια $x$ ισχύει $f(x) > g(x)$; \\
        \textit{Απάντηση:} \dots\dots\dots\dots\dots\dots\dots\dots
\end{enumerate}

\newpage

% ----------------------------------------------------------------------
% 3. ΑΞΙΟΛΟΓΗΣΗ (Exit Ticket)
% ----------------------------------------------------------------------
\begin{center}
  \Large \textbf{Exit Ticket: Τι μάθαμε σήμερα για τη $C_f$;}
\end{center}

\vspace{0.8cm}

\textbf{Ονοματεπώνυμο:} \dots\dots\dots\dots\dots\dots\dots\dots\dots\dots\dots\dots

\vspace{0.8cm}

\textbf{1. Εννοιολογικός έλεγχος:} \\
Αν μια συνάρτηση $f$ έχει πεδίο ορισμού το $[1, 5]$ και τύπο $f(x)=x+2$, ποιο από τα παρακάτω σημεία \textbf{δεν} θα μπορούσε να ανήκει στη γραφική της παράσταση και \textbf{γιατί};
\begin{itemize}
  \item $A(2, 4)$
  \item $B(0, 2)$
  \item $\Gamma(4, 7)$
\end{itemize}
\textit{Απάντηση:} \dots\dots\dots\dots\dots\dots\dots\dots\dots\dots\dots\dots\dots\dots\dots\dots\dots\dots\dots\dots

\vspace{1cm}

\textbf{2. Κατανόηση Γραφήματος:} \\
Συμπληρώστε τις προτάσεις χρησιμοποιώντας τις λέξεις: \textit{τομής, πάνω από, κάτω από, τεταγμένη}.
\begin{itemize}
  \item Η λύση της εξίσωσης $f(x)=0$ αντιστοιχεί στα σημεία \dots\dots\dots\dots\dots\dots με τον άξονα $x'x$.
  \item Η ανίσωση $f(x) < 0$ αληθεύει για εκείνα τα $x$ που η $C_f$ βρίσκεται \dots\dots\dots\dots\dots\dots από τον άξονα $x'x$.
  \item Σε κάθε σημείο $M(x,y)$ της γραφικής παράστασης, η τιμή της συνάρτησης είναι η \dots\dots\dots\dots\dots\dots του σημείου.
\end{itemize}

\vspace{1cm}

\textbf{3. Αυτοαξιολόγηση:} \\
Μετά το σημερινό μάθημα, νιώθω ότι μπορώ να καταλάβω τις ιδιότητες μιας συνάρτησης κοιτάζοντας μόνο το γράφημά της, χωρίς να χρειάζεται να κάνω πράξεις: \\
(Καθόλου) 1 --- 2 --- 3 --- 4 --- 5 (Πολύ)
Ο καθηγητής μου βοήθησε να κατανοήσω τις έννοιες της γραφικής παράστασης μιας συνάρτησης \\
(Καθόλου) 1 --- 2 --- 3 --- 4 --- 5 (Πολύ)

\end{document}